\part*{Part II -- Guide for the student}
\addcontentsline{toc}{part}{Part II -- Guide for the student}
\section{How to use this guide}\label{sec:guide}
Part~I describes \emph{what} the project should find out and \emph{why}. This part is about
\emph{how}: what to do first, which skills to build, how to keep records that will still make sense
in month eleven, and which mistakes cost other people months. Treat it as a starting point and
improve it. If something here turns out to be wrong for our lab, write that down. Corrections like
that belong in your thesis.

\begin{keybox}[Three habits that decide the outcome]
\begin{enumerate}[leftmargin=1.2em]
  \item \textbf{Look before you touch.} Image every end face \emph{before} cleaning or polishing it. A
        cleaned-off film is evidence you can never get back.
  \item \textbf{Log every intervention.} Every re-peak, swap, cleaning, gas-cylinder change and
        power cut goes into the event log with a timestamp. Weeks-long runs are only interpretable
        if you know what happened to them.
  \item \textbf{Write as you go.} Draft the methods chapter while you build. By month~9 you will
        have forgotten why you chose a \SI{5}{mm} lens or a \SI{3}{nm} filter.
\end{enumerate}
\end{keybox}

\section{Your first four weeks}
\begin{table}[htbp]
\centering\small
\caption{Suggested onboarding plan. Adapt it with your supervisor in the first meeting.}
\label{tab:onboarding}
\begin{tabularx}{\textwidth}{@{}l Y L{4.0cm}@{}}
\toprule
Week & Activities & Output \\
\midrule
1 & Departmental, laser, gas and chemical inductions; accounts (lab notebook, repository, data storage);
    read the five starred papers in \cref{sec:reading}; watch an experienced group member polish and
    inspect a connector & reading notes (one paragraph per paper); induction certificates \\
2 & Practise cleaning, inspection and coupling with a visible laser and a spare fibre; measure
    coupling efficiency properly (\cref{sec:skills}); collect degraded fibres and their histories
    from the group & first coupling measurements; inventory of used fibres \\
3 & Phase~0 forensics on the used fibres (imaging, low-power throughput); draft risk assessments;
    first CAD model and a test print with an O-ring groove; Python logger reading one photodiode &
    image set; draft risk assessments; logger prototype \\
4 & Write the project plan (2--3 pages: aims, first experiments, parts list with prices and
    lead times); present it in the group meeting (10 minutes); place orders &
    MS1: plan and risk assessments signed off, orders placed \\
\bottomrule
\end{tabularx}
\end{table}

\section{Core skills}\label{sec:skills}

\subsection{Handling, cleaning and inspecting fibre end faces}
\begin{itemize}
  \item Never touch an end face, and cap every connector as soon as it is unplugged. Store fibres with
        caps on, in a clean box.
  \item Always \emph{inspect first}. If the face is clean, do not clean it. Cleaning tools can leave
        residues and scratch the face.
  \item For cleaning, go from dry to wet: filtered gas or a dry cleaning cassette first, then a
        lint-free wipe moistened with reagent-grade isopropanol, followed immediately by a dry wipe.
        Inspect again.
  \item Use the same microscope, magnification, illumination and focus for every image and save the
        raw files with the fibre ID, end (A/B), date and step in the file name. Use the core and cladding
        zones of IEC~61300-3-35 as the vocabulary for describing defects.
\end{itemize}

\subsection{Polishing}
Learn the group's procedure from an experienced member and then write it down as an SOP with all
parameters (films, grit sequence, times, pressure, lubricant, rinsing). Keep a polishing log for every
end. Variability between polishes is one of the largest uncontrolled factors in this project. Make
polishing boring and reproducible, and measure its outcome (inspection images, coupling efficiency,
end-face geometry if an interferometer is available).

\subsection{Coupling near-UV light into a single-mode fibre}
\begin{itemize}
  \item \textbf{Mode matching.} For a collimated Gaussian beam of $1/e^2$ radius $w_\text{in}$, a lens of
        focal length $f$ produces a waist $w_0=\lambda f/(\pi w_\text{in})$. Choose $f$ so that $w_0$
        equals the fibre mode-field radius. With a mode-field diameter of \SI{\approx2.4}{\micro m} at
        \SI{369.5}{nm}, an $f=\SI{5}{mm}$ asphere needs a \SI{\approx1.0}{mm} diameter beam
        (\cref{app:estimates}).
  \item \textbf{Alignment.} Start with a visible alignment beam if the optics allow it. Make the
        output of the fibre visible, then walk the beam with two mirrors (translation and angle) before
        touching focus. Optimise the focus last and re-walk.
  \item \textbf{PM fibre.} Align the input polarisation to the slow axis: rotate the input half-wave
        plate while measuring the output polarisation extinction ratio, and confirm by gently warming or
        stressing the fibre (the output polarisation should not change).
  \item \textbf{Measure coupling efficiency correctly.} Measure the power directly before the fibre
        (after all optics, including the cell window) and directly after the output, with the same
        calibrated sensor. Correct for the Fresnel loss at each face if you report an ``internal''
        efficiency, and state clearly which definition you use.
\end{itemize}

\subsection{Measuring UV power reliably}
\begin{itemize}
  \item Silicon photodiodes slowly lose responsivity under UV exposure, and photodiodes and amplifiers
        drift with temperature. Keep the intensity on detectors low, prefer integrating spheres or
        diffusers, and consider photodiodes specified for UV resistance for the long runs. (SiC photodiodes,
        often recommended for UV, barely respond at \SI{369.5}{nm}.)
  \item Check the zero (laser off) and the linearity (calibrated attenuator) at the start of every run.
  \item Avoid interference fringes: use wedged optics and tilt surfaces slightly, otherwise
        etalons turn tiny temperature drifts into apparent power changes.
  \item Cross-check all detectors weekly against one reference power meter that is otherwise kept in a
        drawer (so it does not age).
\end{itemize}

\subsection{Gas handling and leak testing}
\begin{itemize}
  \item Get trained on regulators and cylinder handling. Use compression fittings correctly (do not
        over-tighten; replace ferrules) and leak-check every new joint with a leak-detection liquid or a
        helium or hydrogen sniffer.
  \item For ppm-level oxygen, the tubing matters. Polyurethane and fluoropolymer tubes are permeable
        to \Oxy; stainless steel, copper or short nylon lines are better (\cref{app:estimates}).
  \item Leak-test every cell before use: pressurise to about \SI{+100}{mbar} with \Ntwo, log pressure
        \emph{and} temperature for 12 hours, and compute the leak rate $Q=V\,\mathrm{d}p/\mathrm{d}t$ after
        correcting for temperature (\cref{app:sops}).
\end{itemize}

\subsection{CAD and 3D printing of gas-tight parts}
\begin{itemize}
  \item FDM prints are porous unless printed with \SI{100}{\percent} infill, several perimeters, slight
        over-extrusion and good layer adhesion. Even then, seal the inside with a thin low-outgassing epoxy
        for the static cells. Resin (SLA) prints are denser but release residual monomer unless fully
        post-cured. Bake parts (at a temperature safely below the glass transition of the polymer) under
        \Ntwo{} purge before use.
  \item Design O-ring grooves to standard proportions (squeeze and gland fill as in a sealing handbook)
        and print the sealing faces flat and smooth, or face them with a thin epoxy layer. Use heat-set
        metal inserts for screws.
  \item Choose the material with \cref{tab:materials} in mind. PETG is a good first choice for FDM. ABS
        (styrene emission, high water permeability) and standard SLA resins (high outgassing) are poor
        choices near UV-illuminated end faces. Use walls of at least \SI{1.1}{mm} and randomised
        seams~\cite{Gordeev2018}.
  \item Print a test coupon with every new material or setting, and keep coupons for the \arm{N2+P}
        outgassing test.
\end{itemize}

\subsection{Data acquisition and analysis}
\begin{itemize}
  \item Write the logger so that it never loses data. Write to a new file every day, flush after each
        write, use UTC timestamps, and save a metadata file with fibre IDs, arm assignment, calibration
        factors and the software version (git commit). Make it restart automatically after a power cut.
  \item Keep raw data read-only. All corrections (detector temperature, flagged intervals) are applied
        by analysis scripts, never by editing data files.
  \item Make every figure in the thesis reproducible from raw data with one script. A notebook is fine
        for exploration, but the final analysis should be a script in the repository.
  \item Learn enough statistics to defend an upper limit: regression with autocorrelated residuals,
        block bootstrap, and mixed-effects models (\cref{sec:analysis}).
\end{itemize}

\section{Working practices}
\begin{description}[style=nextline,leftmargin=1.2em,itemsep=3pt]
  \item[Lab notebook] Use the group's electronic notebook. Create one entry per day with the plan,
        what was done, observations, and links to data and images. Note surprises immediately, even if
        they seem unimportant.
  \item[Repository] Keep code, CAD files, SOPs, and the thesis source in version control. Commit small
        and often with meaningful messages. Store large raw data on the group's storage, not in git, and
        put the paths in the metadata.
  \item[Meetings] Meet the daily supervisor weekly for 30 minutes with a short written agenda (done,
        problems, plan, decisions needed). Meet the PI monthly, together with a one-page progress
        report that also reviews the risk register (\cref{tab:risks}). Present in the group meeting
        about every two months.
  \item[Asking for help] If you are stuck for more than a day on something that someone in the group
        has done before, ask. Experienced colleagues can often solve fibre-coupling and vacuum or gas
        problems in minutes.
\end{description}

\section{Common pitfalls}
\begin{pitfallbox}[Lessons from long-term optical measurements]
\begin{itemize}[leftmargin=1em]
  \item \textbf{Confusing alignment drift with degradation.} Always record throughput before
        \emph{and} after re-peaking.
  \item \textbf{Cleaning away the evidence.} Image and, if possible, analyse before cleaning.
  \item \textbf{Comparing runs done at different times.} Run conditions in parallel; if you must
        compare runs, keep a reference arm common to both.
  \item \textbf{Trusting a gas cell you have not measured.} A cell with an indicator that has turned
        colour, or an \Oxy{} reading that drifts up, means the run is compromised. Log \Oxy{}, RH and
        flow continuously.
  \item \textbf{Introducing contamination with the fix.} Silicone grease, fresh resin prints, adhesive tapes
        and cable jackets near the end face can all release organics. Keep them out of the cell or
        test them (\arm{N2+P}).
  \item \textbf{Changing two things at once.} New fibre batch, new polish and new atmosphere in the same
        run cannot be untangled.
  \item \textbf{Leaving the thesis until the end.} Writing always takes longer than planned, and writing
        early also shows where data are missing while you can still take them.
\end{itemize}
\end{pitfallbox}

\section{Writing the thesis}
A suggested structure (adapt it to the departmental regulations):
\begin{enumerate}
  \item Introduction: UV light delivery for trapped ions, the problem, aims.
  \item Background: optical fibres and coupling; defects and solarisation in silica; surface
        photochemistry and laser-induced contamination; gas transport and controlled atmospheres.
  \item Methods: testbed, cells and their qualification, procedures, data analysis.
  \item Results I: forensics of used fibres and the baseline degradation law.
  \item Results II: the atmosphere campaigns and the cleaning sequence; surface analysis.
  \item Results III: demonstration on a \SI{369.5}{nm} link.
  \item Discussion: mechanism, comparison with the literature, design guidelines.
  \item Conclusions and outlook.
\end{enumerate}
Appendices: CAD drawings and bill of materials, SOPs, logger documentation, data-availability statement.
Plan for at least eight weeks of focused writing, and give your supervisor complete chapter drafts
early enough for a round of feedback.
